\section{Operation}

This chapter describes the standard operations for running automated image
collection. The goal of normal operation is to define an imaging schedule,
allow the system to automatically acquire images, and prepare the resulting
image dataset for downstream analysis.

During routine experiments, users do not directly interact with the camera.
Instead, experiments are controlled through a schedule file that defines when
images should be captured. The Phenopi scheduler runs continuously in the
background and executes the required captures automatically.

\subsection{Creating an experiment schedule}

Phenopi experiments are defined using a schedule configuration file. This file
describes the duration of the experiment and the times at which images should be
captured.

An example schedule configuration is shown below.

\begin{verbatim}
{
    "start_date": "2026-05-01",
    "num_days": 14,
    "capture_times": [
        "08:00",
        "08:30",
        "09:00",
        "09:30",
        "10:00"
    ]
}
\end{verbatim}

The \texttt{start\_date} field defines the first day of imaging. The
\texttt{num\_days} field determines the total experiment duration. The
\texttt{capture\_times} field specifies the time points at which images are
acquired each day.

Capture times should normally be restricted to the photoperiod of the plants.
For example, when plants are grown under a 08:00--20:00 light cycle, images
should only be scheduled during this period to avoid interrupting the dark
period.

\subsection{Generating the acquisition schedule}

Before an experiment starts, the schedule builder converts the experiment
configuration into a complete list of individual capture jobs.

Each scheduled capture receives an exact timestamp. This generated schedule
represents the complete experimental acquisition protocol and remains fixed
during the experiment.

This approach ensures that image acquisition is reproducible and that the exact
planned imaging times are known before the experiment begins.

\subsection{Running an experiment}

Once the acquisition service is active and a valid schedule has been created,
no further user interaction is required.

The scheduler continuously checks the current time against the experiment
schedule. When a capture time is reached, the scheduler starts a separate
camera process, acquires the image, saves the result, and records that the job
has been completed.

Each capture is independent. The camera is initialized, used for one image, and
then closed again. This reduces the risk of camera instability during
experiments lasting multiple days or weeks.

The scheduler keeps track of completed captures using a persistent state file.
If the Raspberry Pi restarts, previously completed images are not repeated and
the experiment continues from the remaining scheduled captures.

\subsection{Captured image output}

All acquired images are stored in the configured capture directory.

Images are saved using the filename format:

\begin{verbatim}
capture_YYYYMMDD_HHMMSS.jpg
\end{verbatim}

For example:

\begin{verbatim}
capture_20260504_153945.jpg
capture_20260504_154000.jpg
capture_20260504_154015.jpg
\end{verbatim}

The timestamp contained in the filename represents the acquisition time and is
used automatically by the downstream analysis pipeline. Users should therefore
avoid renaming captured images.

Multiple images can be acquired at each time point. These technical replicate
captures allow measurement repeatability to be quantified and help distinguish
true biological changes from image acquisition noise.

\subsection{Monitoring an active experiment}

During an experiment, the scheduler log can be inspected to verify normal
operation.

The service log can be viewed using:

\begin{verbatim}
journalctl -u phenopi-scheduler -f
\end{verbatim}

The log reports scheduled captures, successful acquisitions, and any errors
encountered during operation.

It is recommended to verify the first few captures of every experiment by
checking that images are correctly exposed, focused, and contain the complete
plant tray.

\subsection{After experiment completion}

After the final scheduled capture has completed, the image directory represents
the raw dataset for the experiment.

The capture directory should be preserved unchanged. All downstream Phenopi
analysis tools operate directly on these timestamped images and write processed
outputs to separate results directories.

Keeping the original images unchanged ensures that analyses can always be
repeated with updated parameters or newer versions of the analysis pipeline.
